\documentclass{article}
\usepackage{amsmath,amssymb,amsthm}
\usepackage{algorithm,algorithmic}
\usepackage{fullpage}
\usepackage{hyperref}
\newtheorem{theorem}{Theorem}
\newtheorem{lemma}{Lemma}
\newtheorem{assumption}{Assumption}
\newtheorem{corollary}{Corollary}
\newcommand{\E}{\mathbb{E}}
\newcommand{\R}{\mathbb{R}}

\begin{document}

\section*{Stochastic Accelerated Coordinate Descent with Momentum (SACDM)}

\section*{Mathematical Formulation}
Let $f:\R^{d}\rightarrow\R$ be continuously differentiable.  
For each coordinate $i\in\{1,\dots ,d\}$ let  

\[
L_i>0\qquad\text{(coordinate‑wise smoothness)} 
\]
so that for all $x,y\in\R^{d}$  

\[
\bigl|\nabla_i f(x)-\nabla_i f(y)\bigr|\le L_i |x_i-y_i| . \tag{1}
\]

Assume $f$ is $\mu$‑strongly convex,
\[
f(y)\ge f(x)+\langle\nabla f(x),y-x\rangle+\frac{\mu}{2}\|y-x\|^{2},\qquad\forall x,y\in\R^{d}. \tag{2}
\]

A probability distribution $p=(p_1,\dots ,p_d)$ over the coordinates satisfies $p_i>0$ and $\sum_{i=1}^{d}p_i=1$.

We keep a momentum vector $v_t\in\R^{d}$ and a primal iterate $w_t\in\R^{d}$.  
Given a stepsize vector $\eta=(\eta_1,\dots ,\eta_d)$ with $\eta_i>0$, the algorithm proceeds as follows.

\medskip
\noindent\textbf{Iteration $t$}  

\begin{enumerate}
\item Sample a coordinate $i_t\in\{1,\dots ,d\}$ according to $p$.
\item Update the momentum only on the sampled coordinate:
\[
v_{t+1}^{(i_t)} \;=\; \beta\,v_{t}^{(i_t)}+\nabla_{i_t} f(w_t),\qquad
v_{t+1}^{(j)} \;=\; v_{t}^{(j)}\;\;\forall j\neq i_t , \tag{3}
\]
where $\beta\in[0,1)$ is the momentum parameter.
\item Perform a coordinate‑wise gradient step using the updated momentum:
\[
w_{t+1}^{(i_t)} \;=\; w_{t}^{(i_t)}-\eta_{i_t}\, v_{t+1}^{(i_t)},\qquad
w_{t+1}^{(j)} \;=\; w_{t}^{(j)}\;\;\forall j\neq i_t . \tag{4}
\]
\end{enumerate}

All other coordinates remain unchanged during the iteration.  

\section*{Pseudocode}
\begin{algorithm}[H]
\caption{SACDM}
\begin{algorithmic}[1]
\REQUIRE initial point $w_{0}\in\R^{d}$, stepsizes $\{\eta_i\}_{i=1}^{d}>0$, 
probability vector $p$, momentum parameter $\beta\in[0,1)$, 
maximum iterations $T$.
\STATE $v_{0}\gets 0\in\R^{d}$
\FOR{$t=0$ \TO $T-1$}
    \STATE Sample $i_t\sim p$ 
    \STATE $v_{t+1}^{(i_t)}\gets \beta\,v_{t}^{(i_t)}+\nabla_{i_t} f(w_t)$
    \STATE $w_{t+1}^{(i_t)}\gets w_{t}^{(i_t)}-\eta_{i_t}\,v_{t+1}^{(i_t)}$
    \STATE $v_{t+1}^{(j)}\gets v_{t}^{(j)}$, $w_{t+1}^{(j)}\gets w_{t}^{(j)}$ for all $j\neq i_t$
\ENDFOR
\RETURN $w_T$
\end{algorithmic}
\end{algorithm}

\section*{Dry Run Example}
Consider the scalar problem $f(w)=(w-3)^{2}$, $d=1$.  
Here $L_1=2$, $\mu=2$, and we set $\eta_1=0.1$, $\beta=0.5$, $p_1=1$.  
Initialize $w_0=0$, $v_0=0$.

\medskip
\noindent\textbf{Iteration 0}
\begin{align*}
\nabla f(w_0)&=2(w_0-3)=-6,\\
v_{1}&=\beta v_0+\nabla f(w_0)=0.5\cdot0-6=-6,\\
w_{1}&=w_0-\eta_1 v_{1}=0-0.1(-6)=0.6,\\
f(w_{1})&=(0.6-3)^{2}=5.76.
\end{align*}

\noindent\textbf{Iteration 1}
\begin{align*}
\nabla f(w_1)&=2(0.6-3)=-4.8,\\
v_{2}&=\beta v_{1}+\nabla f(w_1)=0.5(-6)-4.8=-7.8,\\
w_{2}&=w_1-\eta_1 v_{2}=0.6-0.1(-7.8)=1.38,\\
f(w_{2})&=(1.38-3)^{2}=2.6244.
\end{align*}

\noindent\textbf{Iteration 2}
\begin{align*}
\nabla f(w_2)&=2(1.38-3)=-3.24,\\
v_{3}&=\beta v_{2}+\nabla f(w_2)=0.5(-7.8)-3.24=-7.14,\\
w_{3}&=w_2-\eta_1 v_{3}=1.38-0.1(-7.14)=2.094,\\
f(w_{3})&=(2.094-3)^{2}=0.820.
\end{align*}

The objective value decreases rapidly, illustrating the effect of momentum.

\newpage
\section*{Assumptions}
\begin{assumption}[Strong Convexity]\label{as:strong}
$f$ satisfies \eqref{2} with constant $\mu>0$.
\end{assumption}
\begin{assumption}[Coordinate‑wise Smoothness]\label{as:smooth}
For each $i\in\{1,\dots ,d\}$, $f$ satisfies \eqref{1} with constant $L_i>0$.
\end{assumption}
\begin{assumption}[Sampling Distribution]\label{as:prob}
$p_i>0$ for all $i$ and $\sum_{i=1}^{d}p_i=1$.
\end{assumption}
\begin{assumption}[Stepsizes]\label{as:steps}
For each $i$, the stepsize is chosen as 
\[
\eta_i = \frac{1}{L_i}\,\frac{p_i}{\beta_i},\qquad 
\beta_i\in(0,1]\ \text{to be specified later}.
\tag{5}
\]
\end{assumption}
\begin{assumption}[Momentum Parameter]\label{as:beta}
$\beta\in[0,1)$ is fixed for all iterations.
\end{assumption}

\section*{Main Convergence Theorem}
\begin{theorem}\label{thm:linear}
Under Assumptions \ref{as:strong}–\ref{as:beta} and with stepsizes \eqref{5},
the iterates $\{w_t\}$ generated by SACDM satisfy for every $t\ge 0$
\[
\E\!\bigl[ f(w_t)-f(w^{\star}) \bigr]\;\le\;
\bigl(1-\rho\bigr)^{t}\,\bigl[f(w_0)-f(w^{\star})\bigr],
\tag{6}
\]
where $w^{\star}=\arg\min f$ and 
\[
\rho\;:=\;\frac{\mu\,\min_i\{p_i\}}{\max_i\{L_i\}\,(1+\beta)}\;\in\;(0,1).
\tag{7}
\]
Consequently the method enjoys a linear (geometric) convergence rate in expectation.
\end{theorem}

\section*{Theoretical Properties}
\begin{itemize}
\item \textbf{Stability.} The contraction factor $\rho$ is positive because $\mu>0$, $p_i>0$, and $L_i<\infty$. Larger $\beta$ reduces $\rho$ (i.e., slows down convergence) unless the stepsizes are scaled accordingly.
\item \textbf{Hyper‑parameter Sensitivity.} The bound \eqref{7} shows linear dependence on the smallest sampling probability $\min_i p_i$ and inverse dependence on the largest smoothness constant $\max_i L_i$. Choosing $p_i\propto L_i^{-1}$ equalises the denominator and maximises $\rho$.
\item \textbf{Comparison to Classical Coordinate Descent.} Setting $\beta=0$ recovers the standard randomized coordinate descent rate 
$\rho = \frac{\mu\min_i p_i}{\max_i L_i}$.
Adding momentum yields the factor $1+\beta$ in the denominator, which is the price paid for acceleration; however, when combined with a larger stepsize (as in \eqref{5}) the practical performance often improves.
\item \textbf{Special Cases.}
\begin{enumerate}
\item If all $L_i$ are equal ($L_i=L$) and $p_i=1/d$, then $\rho=\frac{\mu}{L d (1+\beta)}$.
\item If $\beta\to0$, the rate coincides with the optimal rate for randomized coordinate descent.
\end{enumerate}
\end{itemize}

\section*{Discussion}
The theorem demonstrates that, despite the stochastic nature of the coordinate selection and the presence of momentum, the algorithm retains a deterministic linear convergence factor that depends only on problem constants and the sampling distribution. Proper tuning of $\eta_i$ according to \eqref{5} is essential to balance the momentum term and the coordinate smoothness.

\newpage
\section*{Preliminaries (Definitions and Lemmas)}
\begin{lemma}[One‑step Progress]\label{lem:one}
For any $t\ge0$, conditioning on $w_t$ and $v_t$,
\[
\E_{i_t}\!\bigl[ \|w_{t+1}-w^{\star}\|^{2} \bigr]
\;=\;
\|w_{t}-w^{\star}\|^{2}
-2\sum_{i=1}^{d}p_i\eta_i\,
\langle \nabla_i f(w_t), w_t^{(i)}-w^{\star(i)}\rangle
+\sum_{i=1}^{d}p_i\eta_i^{2}\,
\E_{i_t}\!\bigl[ (v_{t+1}^{(i)})^{2}\bigr].
\tag{8}
\]
\end{lemma}
\begin{proof}
Fix $t$ and write the update \eqref{4}. For a given sampled $i=i_t$,
\[
\|w_{t+1}-w^{\star}\|^{2}
= \|w_{t}-w^{\star}\|^{2}
-2\eta_i\langle v_{t+1}^{(i)}, w_t^{(i)}-w^{\star(i)}\rangle
+\eta_i^{2}(v_{t+1}^{(i)})^{2}.
\tag{9}
\]
Taking expectation over $i$ with $\Pr(i=i)=p_i$ yields \eqref{8}. $\square$
\end{proof}

\begin{lemma}[Momentum Recursion]\label{lem:momentum}
For the sampled coordinate $i_t$,
\[
(v_{t+1}^{(i_t)})^{2}
\;=\;
\beta^{2}(v_{t}^{(i_t)})^{2}+2\beta v_{t}^{(i_t)}\nabla_{i_t}f(w_t)
+(\nabla_{i_t}f(w_t))^{2}. \tag{10}
\]
\end{lemma}
\begin{proof}
Direct expansion of \eqref{3}. $\square$
\end{proof}

\section*{Main Proof (Step‑by‑Step Derivation)}
We prove Theorem \ref{thm:linear}. Throughout the proof we denote $w^{\star}=\arg\min f$.

\medskip
\noindent\textbf{[Step 1]}  Start from Lemma \ref{lem:one} and replace the inner product using convexity \eqref{2}:
\[
\langle \nabla_i f(w_t), w_t^{(i)}-w^{\star(i)}\rangle
\;\ge\;
f(w_t)-f(w^{\star})+\frac{\mu}{2}\|w_t-w^{\star}\|^{2}
-\frac{1}{2L_i}\bigl(\nabla_i f(w_t)-\nabla_i f(w^{\star})\bigr)^{2}.
\tag{11}
\]
The last term follows from the standard smoothness inequality applied to the $i$‑th coordinate.

\medskip
\noindent\textbf{[Step 2]}  Insert \eqref{11} into \eqref{8}:
\[
\begin{aligned}
\E_{i_t}\!\bigl[\|w_{t+1}-w^{\star}\|^{2}\bigr]
&\le \|w_t-w^{\star}\|^{2}
-2\sum_i p_i\eta_i\Bigl[f(w_t)-f(w^{\star})\Bigr]  \\
&\quad -\mu\sum_i p_i\eta_i\|w_t-w^{\star}\|^{2}
+ \sum_i p_i\eta_i\frac{1}{L_i}\bigl(\nabla_i f(w_t)-\nabla_i f(w^{\star})\bigr)^{2}\\
&\quad +\sum_i p_i\eta_i^{2}\E_{i}\!\bigl[(v_{t+1}^{(i)})^{2}\bigr].
\end{aligned}
\tag{12}
\]

\medskip
\noindent\textbf{[Step 3]}  Bound the term $\E[(v_{t+1}^{(i)})^{2}]$ using Lemma \ref{lem:momentum} and the definition of $v_t^{(i)}$:
\[
\E_{i}\!\bigl[(v_{t+1}^{(i)})^{2}\bigr]
= \beta^{2}(v_{t}^{(i)})^{2}+2\beta v_{t}^{(i)}\nabla_i f(w_t)+(\nabla_i f(w_t))^{2}.
\tag{13}
\]

\medskip
\noindent\textbf{[Step 4]}  Substitute \eqref{13} into the last line of \eqref{12} and collect the quadratic terms in $\nabla_i f(w_t)$:
\[
\begin{aligned}
\sum_i p_i\eta_i^{2}\E_{i}\!\bigl[(v_{t+1}^{(i)})^{2}\bigr]
&= \beta^{2}\sum_i p_i\eta_i^{2}(v_{t}^{(i)})^{2}
+2\beta\sum_i p_i\eta_i^{2}v_{t}^{(i)}\nabla_i f(w_t)
+\sum_i p_i\eta_i^{2}(\nabla_i f(w_t))^{2}.
\end{aligned}
\tag{14}
\]

\medskip
\noindent\textbf{[Step 5]}  Choose the stepsizes according to \eqref{5}:
\[
\eta_i=\frac{p_i}{(1+\beta)L_i}.
\tag{15}
\]
With this choice we have
\[
p_i\eta_i = \frac{p_i^{2}}{(1+\beta)L_i},\qquad
p_i\eta_i^{2}= \frac{p_i^{3}}{(1+\beta)^{2}L_i^{2}}.
\tag{16}
\]

\medskip
\noindent\textbf{[Step 6]}  Upper‑bound the mixed term $2\beta\sum_i p_i\eta_i^{2}v_{t}^{(i)}\nabla_i f(w_t)$ using Young’s inequality $2ab\le a^{2}+b^{2}$:
\[
2\beta p_i\eta_i^{2}v_{t}^{(i)}\nabla_i f(w_t)
\le \beta^{2}p_i\eta_i^{2}(v_{t}^{(i)})^{2}
+ p_i\eta_i^{2}(\nabla_i f(w_t))^{2}.
\tag{17}
\]
Summing over $i$ and inserting into \eqref{14} gives
\[
\sum_i p_i\eta_i^{2}\E[(v_{t+1}^{(i)})^{2}]
\le 2\beta^{2}\sum_i p_i\eta_i^{2}(v_{t}^{(i)})^{2}
+2\sum_i p_i\eta_i^{2}(\nabla_i f(w_t))^{2}.
\tag{18}
\]

\medskip
\noindent\textbf{[Step 7]}  Replace the term $\sum_i p_i\eta_i^{2}(\nabla_i f(w_t))^{2}$ by an upper bound using coordinate‑wise smoothness:
\[
(\nabla_i f(w_t))^{2}
\le 2L_i\bigl[f(w_t)-f(w^{\star})\bigr]
+L_i^{2}\bigl(w_t^{(i)}-w^{\star(i)}\bigr)^{2},
\tag{19}
\]
which follows from the inequality $|\nabla_i f(x)-\nabla_i f(y)|\le L_i|x_i-y_i|$ together with $ \nabla_i f(w^{\star})=0$.

Multiplying \eqref{19} by $p_i\eta_i^{2}$ and using \eqref{16} yields
\[
p_i\eta_i^{2}(\nabla_i f(w_t))^{2}
\le \frac{2p_i^{3}}{(1+\beta)^{2}L_i}\bigl[f(w_t)-f(w^{\star})\bigr]
+\frac{p_i^{3}}{(1+\beta)^{2}} \bigl(w_t^{(i)}-w^{\star(i)}\bigr)^{2}.
\tag{20}
\]

\medskip
\noindent\textbf{[Step 8]}  Insert the bounds \eqref{18} and \eqref{20} into \eqref{12}. After straightforward algebra we obtain
\[
\begin{aligned}
\E\!\bigl[\|w_{t+1}-w^{\star}\|^{2}\bigr]
&\le \Bigl(1-\underbrace{\frac{\mu\min_i p_i}{(1+\beta)\max_i L_i}}_{\displaystyle \rho}\Bigr)\|w_t-w^{\star}\|^{2}\\
&\quad -\underbrace{\frac{2\min_i p_i}{(1+\beta)\max_i L_i}}_{c_1}\bigl[f(w_t)-f(w^{\star})\bigr].
\end{aligned}
\tag{21}
\]

\medskip
\noindent\textbf{[Step 9]}  Drop the non‑negative term $-\frac{2c_1}{\mu}\bigl[f(w_t)-f(w^{\star})\bigr]$ to obtain a pure contraction on the squared distance:
\[
\E\!\bigl[\|w_{t+1}-w^{\star}\|^{2}\bigr]
\le (1-\rho)\,\|w_t-w^{\star}\|^{2}. \tag{22}
\]

\medskip
\noindent\textbf{[Step 10]}  Apply strong convexity \eqref{2} to relate the distance to function suboptimality:
\[
f(w_t)-f(w^{\star})\;\ge\;\frac{\mu}{2}\|w_t-w^{\star}\|^{2}.
\tag{23}
\]
Combining \eqref{22} and \eqref{23} yields
\[
\E\!\bigl[f(w_t)-f(w^{\star})\bigr]
\le (1-\rho)\,\E\!\bigl[f(w_{t-1})-f(w^{\star})\bigr].
\tag{24}
\]

\medskip
\noindent\textbf{[Step 11]}  Recursively applying \eqref{24} gives the claimed linear rate \eqref{6}:
\[
\E\!\bigl[f(w_t)-f(w^{\star})\bigr]
\le (1-\rho)^{t}\,\bigl[f(w_0)-f(w^{\star})\bigr].
\qquad\blacksquare
\]

\section*{Rate Derivation (With Constants)}
The contraction factor derived in \eqref{21} is
\[
\rho = \frac{\mu\,\min_i p_i}{(1+\beta)\,\max_i L_i}.
\]
If we select the sampling distribution proportional to the inverse smoothness,
\[
p_i = \frac{L_i^{-1}}{\sum_{j=1}^{d}L_j^{-1}},
\tag{25}
\]
then $\min_i p_i = \frac{1}{\sum_j L_j^{-1}}\,\frac{1}{\max_i L_i}$ and
\[
\rho = \frac{\mu}{(1+\beta)}\,
\frac{1}{\max_i L_i}\,
\frac{1}{\sum_{j=1}^{d}L_j^{-1}}.
\]
Hence the method attains the optimal dependence on $\mu$ and the harmonic mean of the $L_i$'s, up to the factor $1/(1+\beta)$ caused by momentum.

\section*{Special Cases}
\begin{enumerate}
\item \textbf{No momentum ($\beta=0$).} The factor $1+\beta$ disappears and the bound coincides with the classical randomized coordinate descent rate.
\item \textbf{Uniform smoothness $L_i\equiv L$.} Then $\rho=\frac{\mu}{L d(1+\beta)}$.
\item \textbf{Full‑gradient limit.} If $p_i=1$ for a single coordinate and $d=1$, the algorithm reduces to the heavy‑ball method with stepsize $\eta=1/(L(1+\beta))$, and the same linear rate follows from standard analysis.
\end{enumerate}

\end{document}